% update date: 2024.09.04
% 有問題or更新需求可email至: wufish@gmail.com
% 也可以到github留言: https://github.com/coldwufish/NYCU-thesis-template

% 載入會用到的套件, 通常不需要修改這邊的資料
\input{covers/load_env.tex} 

% 預設英文目錄
% 如果想用中文目錄, 再把下面的註解拿掉
% 【本論文正文為英文, 故維持英文目錄】
%\toggletrue{toc-use-cn}


% #### 下面是自己的資料 ####

% 博士班就把下方註解拿掉吧!!
% 【碩士論文, 保持註解狀態】
%\toggletrue{iamphd}

% 論文名稱
\newcommand{\chineseTitle}{Compile Once, Execute Many：重複性模板實例化任務的可驗證規格編譯方法}
\newcommand{\englishTitle}{Compile Once, Execute Many: Verifiable Specification Compilation for Recurrent Template-Instantiation Tasks}
% 論文名稱太長的話, 可以修改covers/inside_var.tex 的顯示語法, 詳情寫在那邊

% 自己的名字 (英文名字有兩種寫法, 一個是姓在後, 一個是放前面還有逗號. )
\newcommand{\studentCnName}{林冠丞}
\newcommand{\studentEnName}{Kuan-Cheng Lin} % 書名頁
\newcommand{\studentEnNameA}{Lin, Kuan-Cheng} % 封面用

% 教授的名字
\newcommand{\advisorCnName}{陳添福}
\newcommand{\advisorEnName}{Tien-Fu Chen}     % 書名頁
\newcommand{\advisorEnNamex}{Chen, Tien-Fu}   % 封面用

% 共同指導教授的名字, 若有再填寫即可. 目前只支援兩位指導教授的欄位
%\newcommand{\advisorCnNameB}{共同指導教授名字}
%\newcommand{\advisorEnNameB}{OOO Lin}     % 書名頁
%\newcommand{\advisorEnNameBx}{Lin, OOO}   % 封面用

% 論文提交的日期，可以寫口試日期 or 畢業日期
% TODO 口試日期確定後改這兩行（目前暫填 deck 時程推得的 2026 年 11 月）
\newcommand{\ThesisDate}{November 2026}
\newcommand{\ThesisDateTW}{中華民國~ 一一五年十一月}


% [書名頁] 可參考教務處的 學位授予名稱 寫法
% https://aa.nycu.edu.tw/aa/ch/app/statisticalmap/list?module=statisticalmap&id=4358

% 詳情請向系所承辦人員確認，以系所提供的資訊為主

% 下面變數跟表格使用一樣的名稱, 後綴的CN與EN表示中英文
\newcommand{\NameofCollegeCN}{資訊學院} % 這個其實沒用到, 不過為了一致還是放進來
\newcommand{\NameofCollegeEN}{College of Computer Science}
% 只有所、沒有大學部 → 用 Institute
\newcommand{\NameofDepartmentInstituteCN}{資訊科學與工程研究所} % 系所名稱, 請看下方[說明1]
\newcommand{\NameofDepartmentInstituteEN}{Institute of Computer Science and Engineering}
% TODO 向系所承辦人員確認學位名稱與 ResearchTopic 的正式寫法（見下方[說明2]，建議也查 NDLTD 上同所前人論文）
\newcommand{\NameofDegree}{Master of Science} % 學位名稱, 常見的是Master of Science, 不同系所採用不同寫法, 記得查學校的資料.
\newcommand{\ResearchTopic}{Computer Science} % 研究領域, 請看下方[說明2]

% [說明1]
% 這邊需要填寫 XX系 or OO所 的名稱
% 如果你的單位是系所合一, 就用 Department. 如: 土木系碩士班請用 Department
% 只有所, 沒有大學部的話, 就用 Institute. 如: 電信所請用Institute

% [說明2]
% 書名頁的最後會有 in OOOO 的內容. 這個目前找不到通用性的寫法. 
% 大家就自行找前人的畢業論文, 看同系所的人是怎麼寫, 就跟他們寫一樣的吧.
% 不過有些系所不需要寫這個東西, ex: 教育所. 
% 不需要的話就把\ResearchTopic整句刪掉

% --------------------------------------------

% [紙本裝訂順序]
% 封面 >> 書名頁 >> 電子檔著作權授權書 >> 紙本延後公開申請書(如果有) >> 審定同意書 >> 論文內容(致謝-附錄) >> 形式確認書 >> 共同作者貢獻度聲明書(彙編著作者)

% 注意: 有簽名的表單不用放到論文電子檔裡面, 上傳圖書館的時候會有相應欄位讓你上傳.


\makeatletter
% 若已定義則重新定義，未定義則直接建立
\renewcommand*\l@figure[2]{\@dottedtocline{1}{1.5em}{2.3em}{#1}{#2}}
\renewcommand*\l@table[2]{\@dottedtocline{1}{1.5em}{2.3em}{#1}{#2}}
\makeatother

\begin{document}

% ====================================================
% 邊界依「博碩士學位論文格式規範」附件 1／附件 2:
%   封面   上下留邊 3 公分（附件 1）
%   書名頁 上下留邊 2 公分（附件 2）
% 樣板原本兩頁共用 3cm, 與附件 2 不符, 已拆成兩組。
\newgeometry{top=3cm,bottom=3cm,left=3cm,right=3cm}

% 1. 第一頁的封面, 系所資訊使用變數代入
\input{covers/front_var.tex}        % 論文的封面

\newgeometry{top=2cm,bottom=2cm,left=3cm,right=3cm}

% 2. 第二頁的書名頁, 系所&日期使用變數代入
% 目前的浮水印剛好可以把校名&相關資訊包在裡面
\input{covers/inside_var.tex}       % 論文的書名頁

\restoregeometry
% ====================================================
% 書名頁起至最後一頁皆須加入浮水印
% 目前語法不知道怎麼設定透明度, 就直接放刷淡的圖片放在頁面中央
\AddToShipoutPicture{
    \put(-5,-15){
        \parbox[b][\paperheight]{\paperwidth}{%
            \vfill
            \centering
            {\includegraphics[scale=0.65]{covers/logo2.pdf}}%
            \vfill
        }
    }
}
% ====================================================

% [補充] 圖書館有說: 授權書(3)&審定書(5)要裝訂於紙本論文中，不用放在PDF電子檔裡面，要記得請口委簽名。

% 3. 論文電子檔著作權授權書: (要裝訂在紙本論文)

% 4. 博士論文指導教授推薦書(碩士論文免附) 由各所自定，可有可無。檔名 phd_recommend.pdf
%\includepdf[pages={1},pagecommand={\thispagestyle{empty}}]{phd_recommend.pdf}

% 5. 學位論文審定同意書(審定書): (要裝訂在紙本論文)

% 6. 致謝 Acknowledgement {Sections/0.1.Acknowledgement}
% !TEX root = ../main.tex
\thispagestyle{empty}

\begin{center}
\Large
\textbf{誌\hspace{1cm}謝}
\end{center}

\vspace{1cm}
\linespread{2}%
\selectfont
\hspace{0.25cm}

% 下面開始寫致謝內容
% TODO 口試後再寫。以不超過一頁為原則，依個人意願自行決定是否撰寫。
（致謝待撰寫。）

\vspace{3cm}
\begin{flushright}
\studentCnName \ 於

國立陽明交通大學\ \NameofDepartmentInstituteCN

\ThesisDateTW
\end{flushright}\newpage
% 誌謝（依個人意願自行決定是否撰寫）：以不超過一頁為原則。

% 頁碼設定
\frontmatter
\pagenumbering{roman} % 小寫羅馬數字頁碼
{\fontfamily{ptm}\selectfont

% 中英文目錄設定放在covers/toc_var.tex 裡面, 主頁看起來比較清爽, 包含下面項目 (這邊的語法應該不太需要調整)
\input{covers/toc_var.tex}

% 7. 中文摘要及關鍵詞 5-7 個 {Sections/0.2.Abstract_chinese}
% 8. 英文摘要及關鍵詞 5-7 個 {Sections/0.3.Abstract}
% 9. 目錄 toc
% 10. 圖片目錄 lof 
% 11. 表格目錄 lot

% 頁碼設定
\mainmatter
\pagenumbering{arabic} % 阿拉伯數字頁碼

% ====================================================
% 12. 論文正文, 可以每個章節一個.tex檔案
% put your statements in the following
% 維持樣板原本的六章結構與檔名。內容骨架來源: docs/thesis/thesis-outline.md，
% 並依 2026-08-15 的兩項決定調整:
%   (1) 主評測載體由 DWG 翻轉為 docx/xlsx，故舊大綱第 10 章「泛化」不再獨立成章，
%       改為第 3 章的分層論述 + 第 5 章的評測本體，DWG 降為案例研究 (5.9)
%   (2) 舊大綱 §7.2-7.4 由 docs/thesis/evaluation-protocol-v1.md 取代
% 舊大綱 11 章 -> 本論文 6 章的對應:
%   1 Introduction        <- 舊 1
%   2 Related Work        <- 舊 2 (background) + 舊 11 (related work)
%   3 System Model        <- 舊 3 (formulation) + 舊 4 (spec IR) + 舊 6 (runtime)
%   4 Methodology         <- 舊 5 (derivation loop), 本論文的核心演算法
%   5 Performance Eval.   <- 舊 7 (protocol) + 舊 8 (results) + 舊 9 (threats)
%   6 Conclusions         <- 舊 9 (適用邊界) + 大綱結語
% 樣板原附的範例章節移到 Sections/_samples/ (表格、圖片、subfigure 語法可參考)
% !TEX root = ../main.tex
\chapter{Introduction}
\label{ch:intro}

% =====================================================================
% 骨架來源：docs/thesis/thesis-outline.md §1
% 寫作順序：情境 → 現行兩條路都不對 → 我們換一個位置放推論 → 貢獻
% 這一章最後寫。等第 3 章把命題釘死之後再回來。
% =====================================================================

\section{Motivation}
\label{sec:intro:motivation}

% TODO 開場情境：一份模板在實務上會被實例化數十到數千次（不同案場／規格／客戶）。
%      人工填寫耗時且易錯，是自動化的高價值目標。
%      主評測載體是 docx／xlsx 商業表單；DWG 工程圖作為最難的載體獨立呈現。

\section{Two Existing Approaches, Both Wrong}
\label{sec:intro:existing}

% TODO 現行做法一：讓 agent 每次做。問題不只是慢與貴，而是每一次都重新賭一次正確性，
%      而且失敗模式是「沉默失敗」——文件看起來對，某個數字錯了。
% TODO 現行做法二：寫成 skill / SOP 文件。這只是把 prompt 存起來，執行期仍由 LLM 解讀，
%      仍然隨機、仍然吃 token，而且這份 skill 的正確性從來沒有被測試過。
% TODO 【related-work v1 審查要吃下來的三分法，放這裡】
%      RPA = deterministic but manual
%      agents = automatic but stochastic
%      ours = automatic derivation + deterministic execution
%      這是全篇最短的一句定位，不要只放在 related work。

\section{Our Position}
\label{sec:intro:position}

% TODO 把 LLM 用在「一次性推導出規格」，把規格編譯成確定性程式，
%      用可執行 oracle 保證這份規格是對的。O(N)·LLM → O(1)·LLM + O(N)·確定性。
% TODO 【命題已收窄，用這一版，不要用大綱 §0 的舊句】
%      本方法的適用範圍限定於 recurrent template-instantiation tasks whose outputs are
%      deterministically derivable from structured source data and explicit operational rules.
%      收窄的理由見 benchmark-case-v1-thought.md §8.6 與 related-work v1 審查結論。
% TODO 【單一命題，避免連言式新穎性】related-work 審查指出五個 gap 都是
%      「別人各自有、只有我合起來」。重心要收到一句：
%      在什麼可執行條件下，允許把 LLM 從執行路徑上永久移除？

\section{Contributions}
\label{sec:intro:contributions}

% TODO C1–C5，見 thesis-outline.md §0 的貢獻表。
%      注意 C4（zero-LLM execution）被 related-work 審查判定為唯一乾淨的 claim，
%      其他四項都有危險鄰居，論述時要放在 C4 之下而不是並列。

\section{Claims and Where They Are Tested}
\label{sec:intro:claims}

% TODO 三個 claim，每一個都要在第 7 章有對應實驗：
%      1. 正確性不輸逐實例 agent，且執行期變異度為零
%         （但推導期變異度不是零，必須誠實報告，見第 7 章）
%      2. 攤提後成本低一到數個數量級，存在明確 break-even 點
%      3. 失敗會「吵」而不是沉默：模板改版時能被定位、能被修復

\section{Thesis Organization}
\label{sec:intro:organization}

% TODO 最後寫。
 \newpage
% !TEX root = ../main.tex
\chapter{Related Work}
\label{ch:relatedwork}

% =====================================================================
% 骨架來源：docs/thesis/thesis-outline.md §2（background）+ §11（related work）
% 兩者合併成樣板原本的第 2 章。前四節是背景與問題刻畫，後七節是文獻。
%
% 相依：docs/thesis/related-work-v1-thought.md（已離線審完）
%
% 【使用該產出前必讀】它的「已驗證」實際只等於「連結打得開」：
%   - 至少 12 篇有同儕審查場所的論文被標成 arXiv preprint（場所欄要用 DBLP 重掃）
%   - 六篇 2026 的論文無法離線核對，其中 arXiv 2607.08010（Amazon tool-making）
%     一篇獨力推翻桶 1 的 gap，核對完成前不得據此改任何內容
%   - 桶 3（LLM + PBE）時間軸停在 2023，2024–2026 全空白，這是唯一非補不可的洞
% =====================================================================

\section{Template-Shaped Tasks}
\label{sec:rel:tasks}

% TODO 任務定義：給定模板 T 與一組值 V，產生實例 D = f(T, V)。
%      難點在 f 的定位（locator）部分：哪一個儲存格／bookmark／文字對應哪一個語意欄位。

\section{Document Formats and Representation Mismatch}
\label{sec:rel:formats}

% TODO docx / xlsx 的資料模型：bookmark、content control、具名範圍、合併儲存格、
%      公式 vs 快取值。以及 DWG 的 entity graph、layer、block/attribute、xdata。
% TODO 表徵不匹配：語意存在於佈局關係，序列化成 token 會摧毀這些結構。
%      【降階處理】這是動機的一部分但不能當主論點，會被一句「下一代模型就解決了」擊破。
%      正確的推論是：用工具提供結構化 dump，並且只在編譯期讀一次。

\section{Failure Modes}
\label{sec:rel:failures}

% TODO 本節是後面 oracle 與 gate 設計的依據，第 3/4 章要有一張對照表回指這裡：
%      - 沉默失敗（值錯、位置對）
%      - 定位漂移（模板改版後 locator 指到別的地方）
%      - 覆蓋不足（欄位根本沒被規格涵蓋 -> 保持模板預設值，最難察覺）
%      - 格式／單位錯誤
%      - negative obligation 違反（該空的沒清空）—— 新增，來自 benchmark S5
%      - 範圍外延錯誤（聚合多算或少算一列，數字有理有據但錯）—— 新增，來自 ADR 0002

\section{Governance Requirements}
\label{sec:rel:governance}

% TODO 可重現、可稽核、可簽核、可回溯。這是純 agent 方案的結構性障礙。
% 【措辭修正】不要寫「法規普遍要求 deterministic output」——找不到學術來源。
%      改寫成 traceability / reproducibility / change control。

\section{LLMs for Document and Engineering-Drawing Understanding}
\label{sec:rel:understanding}

\section{Document Automation, RPA, and Template Filling}
\label{sec:rel:rpa}

% TODO 三分法在這裡展開，但結論句已經在 Introduction 出現過：
%      RPA = deterministic but manual / agents = automatic but stochastic /
%      ours = automatic derivation + deterministic execution

\section{Programming by Example and Program Synthesis}
\label{sec:rel:pbe}

% TODO 【這一節是第二危險的鄰居，也是產出裡唯一非補不可的洞：2024–2026 全空白】
%      FlashFill++ 一類的工作是最直接的比較對象。
%      區別點：範圍聚合（S3）與 negative obligation（S5）是兩個無法化約成單欄 PBE 的案例。

\section{Tool-Making, Skill Libraries, and Amortized LLM Computation}
\label{sec:rel:toolmaking}

% TODO 【第一名危險鄰居】LATM、Voyager skill library、TroVE、CRAFT、Agent Skills，
%      以及待核對的 arXiv 2607.08010（Amazon tool-making）。
%      Gap：這些工作把產生的工具存下來就結束——沒有可執行 oracle、沒有覆蓋率概念、
%      沒有回歸與版本治理。而在有責任的交付物上，這三者才決定能否上線。

\section{Execution-Feedback Program Repair}
\label{sec:rel:repair}

% TODO Self-Debug、Reflexion、AlphaCodium。
%      Gap：它們的 oracle 多來自現成單元測試；本工作的 oracle 必須從文件比對建構，
%      而且要處理「沉默失敗」而非「程式當掉」。

\section{Verification and Guardrails for LLM Output}
\label{sec:rel:verification}

% TODO LLM-as-judge 的侷限、property-based / metamorphic testing。
%      支撐第 4 章的「要 oracle 不要 judge」與第 5 章評測器不得問 LLM 的決定。

\section{Cost-Aware Agent Evaluation}
\label{sec:rel:evaluation}

% TODO Gap：多數 agent 評測只報準確率，不報變異度與攤提成本，
%      而這兩項正是生產環境的決策依據。這一節直接支撐第 5 章的指標選擇。
  \newpage
% !TEX root = ../main.tex
\chapter{System Model}
\label{ch:architecture}

% =====================================================================
% 骨架來源：docs/thesis/thesis-outline.md §3 + §4 + §6，併入舊 §10（泛化）的分層論述
%
% 【本章與第 4 章的分界】
%   第 3 章 = 靜態存在的東西：問題的形式化、規格長什麼樣、規格怎麼被執行
%   第 4 章 = 它怎麼被產生出來的：derive -> gate -> oracle -> repair
%
% 【結構調整說明】舊大綱把「泛化到 docx/xlsx」放成獨立的第 10 章最小 PoC。
% deck 已把主評測載體翻轉成 docx/xlsx，所以泛化不再是附加章節，
% 而是本章 §3.2 的分層論述 + 第 5 章的評測本體。DWG 降為最難的載體。
% =====================================================================

\section{Problem Formulation}
\label{sec:design:formulation}

\subsection{A Stateful, Two-Stage Input Model}
\label{sec:design:inputs}

The abbreviation $\mathcal{A}:T\rightarrow S$ is potentially misleading because the
system does not consume all experimental material in one call, and derivation is neither
pure nor deterministic. Let $d$ denote one derivation dataset in template lineage
$\ell$. We split the available material into a derivation layer and a later regression
layer,
\begin{align}
T &= \bigl(T_{\mathrm{der},d}^{(n)},T_{\mathrm{reg},\ell}\bigr), \\
T_{\mathrm{der},d}^{(n)}
  &= \bigl(d,\ell,B_d,Y_d^{?},\{X_{dj}\}_{j=1}^{N_d},R_d^{?},U_d^{?},
      Q_d,P,\Gamma,\Theta,H_d^{(<n)}\bigr), \\
T_{\mathrm{reg},\ell}
  &= \mathcal{R}_{\ell}
   = \left\{c_i=\bigl(I_i,\{X_{ij}\}_{j=1}^{N_i},u_i,\pi_i,Y_i^{?}\bigr)
     \right\}_{i=1}^{K_{\ell}} .
\end{align}
Here, $B_d$ is exactly one baseline template; $Y_d^{?}$ is an optional expected
output; $\{X_{dj}\}_{j=1}^{N_d}$ is a possibly empty collection of source artifacts;
$R_d^{?}$ is an optional operational-rule artifact; and $U_d^{?}$ is an optional
inline source-data override. When present, the override bypasses the locally parsed
artifact dictionary and, on the agent path, takes precedence over extracted keys. $Q_d$
groups the inspection level, force-refresh choice,
content-addressed dump, cache state, CAD preflight and health observations, locator
verification, and deterministic validation/repair feedback. The derivation environment
is also an input: $P$ denotes the prompt files, $\Gamma$ the agent's tool and skill
definitions, and $\Theta$ the model, gateway, and MCP configuration. Thus, changing a
prompt, an allowed tool, or the selected model changes $T_{\mathrm{der}}$ even when all
customer artifacts are byte-identical.

The term $H_d^{(<n)}$ is the pre-existing state before the $n$th derive. It contains
the field names and rule inventories accumulated from all existing non-degraded specs
derived from the same dataset. The implementation reads this state through
\texttt{derivation\_service.DerivationService.\_reference\_specs},
\texttt{\_field\_hints}, and \texttt{\_previous\_rule\_inventory}; the degraded flag
used by that selection is produced by \texttt{cad\_health.degradation}. The hints are
fed to the agent and the prior inventory is merged into the new spec evidence. With
$\omega_n$ denoting model sampling and the observed session/tool trace, the honest
derivation relation is therefore
\begin{equation}
S_d^{(n)} \sim
\mathcal{A}\!\left(T_{\mathrm{der},d}^{(n)};\omega_n\right),
\qquad
T_{\mathrm{reg},\ell}\notin\operatorname{dom}(\mathcal{A}).
\label{eq:stateful-derivation}
\end{equation}
Consequently, two derives over the same artifact bytes need not have the same context or
the same output: $H_d^{(<n)}$ may have changed, and $\omega_n$ is stochastic. In this
thesis, $\mathcal{A}:T\rightarrow S$ is only shorthand for the projection
$\mathcal{A}:T_{\mathrm{der}}\rightsquigarrow S$ in
Equation~\ref{eq:stateful-derivation}, not a claim that $\mathcal{A}$ is a pure function.

This split follows the implemented call boundary. The dataset creation and derive
paths, \texttt{api.create\_dataset} and
\texttt{derivation\_service.DerivationService.derive}, admit one baseline, at most one
expected artifact, $N_d$ source artifacts, and at most one rules artifact. They do not
read \texttt{models.DwgTestCase}. Only after $S$ exists does the separate endpoint call
\texttt{regression.RegressionRunner.run}, which selects every active test case of the
spec's lineage and executes them. $K_{\ell}$ is therefore the number of active cases at
regression time; the implementation imposes no $K_{\ell}=3$ limit. Each case $c_i$
contains an input template $I_i$, its own source artifacts and explicit source-data
override $u_i$, runtime parameters $\pi_i$, and an optional expected output $Y_i^{?}$.

The prompt, tools, and model terms above are concrete inputs rather than abstract
nuisance variables. \texttt{derivation\_agent.OpenCodeAgent.derive\_rules} selects the
\texttt{dwg-derivation} agent and \texttt{OpenCodeAgent.\_run} submits the task payload;
the selected agent is configured by \texttt{opencode.jsonc}, its prompt is
\texttt{prompts/derivation.md}, and its permitted CAD MCP tools are defined by the agent
configuration and \texttt{mcp\_server.server}. These files and settings must therefore
be versioned with the material definition used in an experiment.

We reason about the resulting specification as
\begin{equation}
S=\bigl(G,L_s,L_t,F,\Pi,\mathcal{E}_S\bigr),
\end{equation}
where $G$ contains value-source and transformation semantics, $L_s$ contains source
locators, $L_t$ contains independently derived target locators, $F$ contains rendering
rules, $\Pi$ is the runtime-parameter schema, and $\mathcal{E}_S$ is derivation and
validation evidence. The persisted \texttt{DwgSpec} stores these concerns across
\texttt{field\_map}, rule projections, \texttt{runtime\_params}, and
\texttt{evidence}; this tuple is a conceptual separation rather than a claim that each
component is a distinct database column.

\paragraph{Information boundary.}
The operational-rule artifact $R_d$ may identify source fields and state formulas,
conditions, constants, and formatting requirements, but it must not contain target cell
addresses, bookmark names, content-control tags, or mappings of the form
\texttt{source\_field $\rightarrow$ target\_cell}. Target references belong to $L_t$ and
must be derived from the template by $\mathcal{A}$; otherwise $R_d$ becomes an answer
table and the derivation problem is leaked into its input. \textbf{Status:
normative/aspirational only.} The boundary is stated but not enforced, in two
independent ways. First, no admission check exists: neither
\texttt{api.create\_dataset} nor
\texttt{derivation\_service.DerivationService.\_load\_rules} rejects an answer-bearing
rules artifact, so a hand-authored $R_d$ containing target addresses is accepted
silently. Second, the persisted projection is not structurally safe:
\texttt{derivation\_service.DerivationService.derive} copies the field's target-reference
slot into \texttt{DwgSpec.mapping\_rules}. That slot is carrier-agnostic --- its
interpretation follows the target kind of the template, so it holds bookmark identities
for \texttt{.docx} and cell addresses for \texttt{.xlsx}, which are exactly the two
artefact classes this boundary forbids. (The field retains the legacy name
\texttt{handles} from the drawing-oriented prototype; the naming is historical and does
not narrow the scope.) If that projection is exported, presented as the rule table, or
fed back as $R_d$, it exposes target answers, and $\mathcal{A}$ is then measured on its
ability to copy them rather than to derive them. It must therefore remain internal, and
be stripped of target references or split from the semantic rules, before the boundary
may be described as a property of the system rather than a requirement on it.

\subsection{Compile-Once, Execute-Many}
\label{sec:design:compile-once}

The target problem class consists of recurrent template-instantiation tasks for which
the output is deterministically derivable from structured source data and explicit
operational rules. ``Compile once'' means that the expensive, stochastic derivation in
Equation~\ref{eq:stateful-derivation} produces one reusable $S$; it does not mean that a
single fixed output job is reused, because instance values differ. For instance input
$V_i=(I_i,\{X_{ij}\},u_i,\pi_i)$, execution deterministically instantiates and applies
the spec,
\begin{align}
(J_i,P_i) &= \mathcal{C}(S,V_i), \\
D_i &= \mathcal{X}(J_i,I_i), \\
\operatorname{ok}_i
  &= \mathcal{O}_{\mathrm{plan}}(P_i,I_i,D_i)
     \land
     \begin{cases}
       \mathcal{O}_{\mathrm{expected}}(D_i,Y_i), & Y_i^{?}=Y_i,\\
       \mathrm{true}, & Y_i^{?}\text{ is absent}.
     \end{cases}
\end{align}
$\mathcal{C}$ is the deterministic per-instance compiler, $P_i$ is its authoritative
write plan, $\mathcal{X}$ is the format-specific executor, and $D_i$ is the produced
document. The problem is to derive an $S$ once such that $\operatorname{ok}_i$ holds
for every admissible instance, while the cost of derivation is amortized over repeated
executions.

The deterministic core is implementation-backed. In
\texttt{compiler.compile\_cad\_job}, the spec, source data, runtime parameters, and
current template values are explicit inputs; failures are returned before submission,
and the canonical job JSON is content-hashed. The DOCX/XLSX path shares this compiler
and applies $P_i$ through \texttt{docx\_target.DocxTargetAdapter.apply} or
\texttt{xlsx\_target.XlsxTargetAdapter.apply};
\texttt{document\_production.document\_job} records a
canonical hash for the document job. Regression is also implementation-backed:
\texttt{regression.RegressionRunner.\_run\_case} compiles each case, and
\texttt{regression.RegressionRunner.\_verify\_case} uses
\texttt{verification.verify\_output} plus
\texttt{verification.compare\_expected} when $Y_i$ exists.

This formulation deliberately localizes randomness to $\mathcal{A}$, and the execution
path now matches that localization. An earlier revision of this chapter recorded that it
did not: \texttt{production.ProductionController.start\_run} invoked
\texttt{derivation\_agent.OpenCodeAgent.repair\_cad\_job} after a failed gateway job
validation, and the comparison that guarded the repaired job excluded target locations,
so a model-relocated target could be accepted and submitted inside the same run. That
branch has since been removed from execution. A compiled job the tenant manifest rejects
terminates the run with \texttt{CAD\_JOB\_MANIFEST\_REJECTED} and an instruction to
re-derive, review, and republish; a cached target reference the current template no
longer contains terminates it with \texttt{CACHED\_HANDLE\_NOT\_FOUND}. Neither is
repaired in place, and \texttt{repair\_cad\_job} has no remaining caller on the execution
path.

\textbf{Status: implementation-backed, and pinned by a negative control rather than by
inspection.} \texttt{test\_successful\_production\_operation\_records\_zero\_llm\_calls}
replaces both agent repair entry points with a recorder that returns a failure, runs a
complete production operation, and then asserts both that the recorder was never invoked
and that the \texttt{LlmOperation} persisted for the run holds zero model calls. Reading
the source would only show that the call site is absent today; the test also fails if a
future change reintroduces one. This is the evidence behind claim C4, and it is
deliberately the narrowest of the thesis's claims: it covers $\mathcal{X}$, not
$\mathcal{A}$, and not the derivation-time repair loop of Chapter~\ref{ch:method}, which
is by construction inside the compile step whose cost repeated execution amortizes.

\subsection{Required Properties of the Specification IR}
\label{sec:design:properties}

The reusable IR must satisfy the following five properties. Each status below refers
to the repository state examined for this thesis, not merely to the intended design.

\paragraph{P1---separation of structure and instance values.}
$S$ must describe reusable source/target locators, transformations, formatting, and
parameter types; a particular instance supplies values through $V_i$. The only values
permitted in $S$ are constants explicitly prescribed by $R_d$, not values copied from a
calibration output. \textbf{Status: implementation-backed at the representation
interface.} \texttt{compiler.compile\_cad\_job} accepts \texttt{spec},
\texttt{source\_data}, and \texttt{runtime\_params} separately, while
\texttt{value\_layer.resolve\_field\_raw} evaluates the value layer at execution time;
\texttt{regression.RegressionRunner.\_run\_case} reuses the same spec with each case's
inputs.
This interface supports reuse, although successful regression---not the type shape
alone---is what detects a semantically overfit constant.

\paragraph{P2---explainable and verifiable locators.}
Every source or target locator must expose what it resolved to and must be checkable
against the corresponding artifact; an unresolved or ambiguous locator is evidence, not
permission to guess. \textbf{Status: implementation-backed.}
\texttt{source\_locator.resolve\_locator} returns the resolved sheet, cell, selection,
and match count, and \texttt{source\_resolution.verify\_against\_sample} checks that the
locator reproduces the derivation sample. On the target side,
\texttt{target\_ref.BaseTargetAdapter.verify} and \texttt{describe} provide a common
existence and explanation interface.
\texttt{docx\_target.DocxTargetAdapter.inspect} and
\texttt{xlsx\_target.XlsxTargetAdapter.inspect} retain labels, physical locations,
positional flags, and ambiguous references, while their \texttt{apply} methods return
every missed reference.

\paragraph{P3---bounded layout resilience and fail-loud behavior.}
Semantic identities should survive admissible layout changes, while a change outside
the declared tolerance must stop the run rather than trigger a guessed write. This
includes the zero-LLM execution requirement: an execution-time model must not relocate a
target. \textbf{Status: implementation-backed for the fail-loud half; the resilience
half is bounded and measured rather than general.} Source locators resolve by
sheet/header/section semantics; \texttt{docx\_target.DocxTargetAdapter} distinguishes
stable content-control/bookmark identities from positional table cells;
\texttt{xlsx\_target.XlsxTargetAdapter} distinguishes stable defined names from
positional coordinates; and both reject duplicate or missing references. The execution
branch that previously undermined this property --- agent repair of a rejected job
inside the run --- has been removed, as described in
Section~\ref{sec:design:compile-once}, so a target that cannot be resolved now ends the
run with a code instead of being relocated.

What is \emph{not} claimed is the resilience half in general form. Layout tolerance is
declared per locator mechanism, and a change outside a mechanism's declared tolerance is
detected only insofar as the mechanism can see it. Two measured boundaries are carried
forward to Section~\ref{sec:design:boundaries} rather than glossed here: an unused
detail-row cell that no column of a repeat block addresses is not cleared by
\texttt{unused\_slots: "blank"}, so a template reused from a previous reporting period
can carry a stale value into a fresh run; and the derivation-time semantic oracle
compares only references whose text differs between the baseline and the expected
document, so a specification that writes into a cell the expected document leaves blank
is structurally invisible to it.
\paragraph{P4---diffability, versioning, and lineage.}
The IR and its evidence must support reviewable changes, immutable publication, and
regression over a template lineage. \textbf{Status: implementation-backed for the
currently produced handle-backed field-map shape.}
\texttt{models.DwgSpec} carries a lineage-scoped unique version and template fingerprint;
\texttt{drift.spec\_field\_drift} reports added, removed, and attribute-level field
changes; \texttt{models.DwgProductionVersion.save} prevents mutation of the published
spec/hash identity; and \texttt{regression.RegressionRunner.run} persists results for
all active cases of the lineage. The drift function reads \texttt{target.handles} rather
than the generic \texttt{targets} representation; diffability for a future spec that
uses only generic target refs remains normative until that path is implemented.

\paragraph{P5---independently auditable rule coverage.}
$S$ must enumerate which rules and fields it claims to cover, and the denominator must
come from the rule artifact independently of the model being scored. \textbf{Status:
implementation-backed since the coverage-v2 path replaced agent self-enumeration.} The
original implementation --- \texttt{rule\_coverage.merge\_inventory} over an agent-supplied
\texttt{rule\_inventory}, unioned across prior derives --- did not satisfy P5 and is
retained only to read historical specifications. A rule the agent never enumerated was
absent from the denominator, so ``22 of 22 covered'' and ``22 of 40 covered'' were the
same observation.

The active path externalizes the denominator instead. A dataset that carries a rule
document must also carry a frozen, versioned rule-applicability manifest;
\texttt{rule\_coverage.parse\_rule\_manifest} verifies that the manifest names the exact
SHA-256 of the rule-document bytes and that its declared document roles match the
dataset's, and \texttt{derivation\_service.DerivationService.\_load\_rule\_manifest}
refuses the derive with \texttt{RULE\_MANIFEST\_MISSING} before any paid model call, so a
run cannot discover after the fact that its release denominator is unauditable. The agent
supplies a \emph{disposition} for each externally assigned rule identifier and nothing
else: \texttt{rule\_coverage.score\_manifest\_coverage} discards unknown identifiers,
refuses duplicate dispositions, and refuses an assignment whose document roles disagree
with the manifest, so the scored agent can neither add, remove, rename, nor move a rule
between roles. The resulting partition into \texttt{covered}, \texttt{unresolved},
\texttt{rejected}, and \texttt{unmentioned} is mutually exclusive and sums to the frozen
applicable set, and \texttt{gates.canonical\_unresolved} turns every non-\texttt{covered}
disposition into a publish gate rather than partial credit. Section~\ref{sec:der:coverage}
develops what this measures and what it still does not.
\subsection{Measured Boundaries of the Specification IR}
\label{sec:design:boundaries}

The five properties above state what the IR must do. This subsection states what it
demonstrably does not, because a boundary that is only discovered by a reader running the
system is indistinguishable from a defect the author did not notice. Each item below is
either pinned by a test or was measured on a real artifact; where an alternative shape was
considered and rejected, the reason is given, since ``we did not think of it'' and ``we
rejected it'' are different claims.

\paragraph{B1---\texttt{exclude} is a row-content predicate.}
A range's \texttt{exclude} clause inspects the content of a data row. It does not inherit
meaning from a preceding section header, so a sheet whose rows alternate between
``current month'' and ``cumulative'' groups under one header cannot be partitioned by
\texttt{exclude}. The rejected alternative was a domain flag such as
\texttt{exclude\_cumulative}; it was refused because it names one agency's form inside a
carrier-neutral IR. The general shape adopted instead is a group projection
(\texttt{records}: a fixed \texttt{group\_size} with named row roles and per-column
\texttt{select\_offset}), which describes the same layout without naming the domain.

\paragraph{B2---stacked records assume a fixed group size.}
\texttt{records} models ``$g$ physical rows form one logical record'' for a constant $g$.
A record whose row count varies with its own data is outside the IR. Both the read side
and the write side are implemented, and a specification produced by the agent for the
XLSX evaluation template declares \texttt{records} with \texttt{group\_size: 2} and named
row roles, so the primitive is exercised by a main evaluation dataset rather than by
fixtures alone.

\paragraph{B3---the verification basis for the collection primitives is $n=1$.}
\texttt{artifact\_collection}, \texttt{records}, and the equality merge of
Section~\ref{sec:design:loc:range} each have exactly one motivating real workflow. Six
independently sourced government workbooks were scanned for a second instance of a
repeating multi-row record block and none was found; the nearest candidate was a one-off
label merge, not a repeating record. The generalization threshold this thesis sets for a
new primitive is therefore not met for these three, and they are presented as
carrier-neutral shapes with a single validating workflow rather than as generally
validated primitives. The risk this creates is specific and worth naming: it does not
appear as a failing test, because the tests were written from the same single example.
It appears as a wrong result on a form nobody has tried.

\paragraph{B4---merge is equality-only and inner.}
Two resolved row sets can be joined on equal key columns. Range containment, interval
overlap, and inequality correspondences are not expressible. The join is deliberately
inner: supplying a default for an unmatched key would convert a missing record into a
normal-looking number without anyone having authorized the default, which is precisely the
silent-failure shape this thesis measures. A specification that needs the filled row must
produce it on the source side, where the decision is visible. Merges chain, so three
tables can be joined as two binary merges; the agent produced a chained pair unprompted.

\paragraph{B5---the aggregation axis is columns only.}
Aggregation sums named columns. The criterion is who decides the number of columns: if the
template does, naming them is legitimate; if \emph{this instance's data} does, the
specification has become instance-specific, and the correct modelling is a long table plus
a pivot at write time. A row axis is not offered because a wide table's column names would
then be data values, which would disable the author-time static check that a formula only
names columns the row set actually has.

\paragraph{B6---a repeat block clears only the cells its own columns address.}
When a block declares \texttt{unused\_slots: "blank"}, the cells cleared in an unused slot
are the cells that block's columns point at. A cell in the same slot that no column
addresses is not cleared, because the IR knows a block's targets only through the strings
its columns declare, and ``every other cell on this row'' is not something it can name.
For a pristine template this is exactly right, and every instance of both evaluation
datasets starts from a pristine template. For the common practice of reusing last period's
completed report as this period's template, it means a stale value can sit on a row the
system considers written. The behavior is pinned by a test rather than left to inference,
and the rejected workaround---declaring a column whose value is always blank---was refused
because it puts layout knowledge into the value layer and makes the column count grow with
the template's width.

\paragraph{B7---one instance is one filled target template.}
A specification may declare several output documents, but their roles must be distinct; a
run that would produce $M$ documents of the same role is refused. Relaxing this is a
separate decision and is not claimed here.

\subsection{Design Principle: Fail Loud, Not Silent}
\label{sec:design:failloud}
% TODO ★ 尚未動筆 ★ 這一小節在 2026-08-15 撰寫 §3.1 時被整段刪除，同日復原為寫作規格。
%      刪除它不是決定，是疏漏——本節是全篇的核心設計信條，不可省略。
% TODO 明確寫成本論文的核心設計信條，並回指第 2 章的失敗模式分類。
%      要點：一個有理有據但錯的輸出，比一個明白失敗的輸出更危險，因為它會被交付出去。
%      §3.1.1 的 information boundary 與 §3.1.3 的 P3、P5 都是這條信條的實例，
%      寫的時候要把它們收攏成同一個原則，不要各講各的。

\subsection{Scope and Applicability Conditions}
\label{sec:design:scope}
% TODO ★ 收窄後的命題寫在這裡，Introduction 只放結論句 ★
%      2026-08-15 註：命題本身已寫進 §3.1.2（Compile-Once, Execute-Many）的正文，
%      但那裡是順著 compile-once 的論證帶出來的，沒有獨立成「適用條件」的判準清單，
%      也沒有往第 6 章的前向引用。本小節仍要寫，並與 §3.1.2 交叉引用而非重複。
%      適用條件（三條都要成立）：
%        1. recurrent —— 同一份模板會被實例化多次（否則過不了 break-even）
%        2. deterministically derivable —— 輸出可由結構化來源資料確定性推導
%        3. explicit operational rules —— 存在明確的操作規則文件，而非需要真實推理的判斷
%      不適用的情境在第 6 章 §6.2 展開，這裡只給定義不談後果。

% =====================================================================

\section{Domain-Dependent and Domain-Independent Layers}
\label{sec:design:layers}

% TODO 這一節決定論文被當成 tool paper 還是 method paper，要放在系統描述的最前面。
%      領域相關：dump 工具、locator 的實作、target adapter
%      領域無關：IR 結構、gates、oracle 種類、repair loop、lineage、覆蓋率計分
%      三條載體（docx / xlsx / dwg）共用同一條 derive -> gate -> compile -> run 管線。

\section{Specification IR}
\label{sec:design:ir}

% TODO 欄位定義、locator、值來源、格式化規則、約束。

\section{Structured Dump as Grounding}
\label{sec:design:dump}

% TODO docx / xlsx 的 inspect 與傾印；CAD 端的 DumpDwg 與 labeled-field 解析。
% TODO 【誠實揭露】openpyxl 的 load/save 有量測到的損失（儲存格註解、列印設定），
%      python-docx 亦有邊界情形。這些屬於載體限制而非方法限制，但要寫出來。

\section{Locator Mechanisms}
\label{sec:design:locators}

\subsection{Semantic Locators}
\label{sec:design:loc:semantic}

% TODO label -> value 的鄰近／對齊關係；以 section 限定搜尋範圍避免跨區塊誤命中；
%      sheet alias 與 fallback 策略。

\subsection{Range Locators and Source-Side Aggregation}
\label{sec:design:loc:range}

% TODO 【寫作規格 2026-08-23 更新——原註記說「實作尚未開始」，那是錯的，
%      而且方向是低估系統。ADR 0002 是 2026-08 以後每一份規格的地基，
%      它的讀取端、顯式邊界、來源聚合與 evidence 都在 production path 上。
%      照舊註記寫會把已完成的能力寫成未來工作。】
%
%      本小節要寫的是**一條演進線**，不是一個機制，因為每一步都是被實測逼出來的：
%        1. ADR 0002（Accepted，G1–G4／G6 已實作；G5 是 benchmark case assertion）
%           範圍定位器 ＋ 來源端聚合。約束：外延必須顯式宣告，不允許隱式整欄。
%           難題是外延：明細第幾列到第幾列、含不含小計列、插列後還對不對。
%           外延差一列會產生「有理有據但錯」的數字，是沉默失敗的教科書範例。
%        2. ADR 0004／0005：把一段列壓成純量不夠——spec 只觸及模板 14 個明細列位的
%           前 3 個，因為校準實例只有 3 筆。那不是 compile-once 是 compile-per-instance。
%           修法是列投影 ＋ repeat 區塊 ＋ 具名外延 ＋ 多文件。
%        3. ADR 0006：repeat 的欄只能直接讀一個來源欄，做不到逐列計算與跨檔查表。
%           修法是**外延先解析成 row set**，逐列公式、join、合計都在那張表上做。
%        4. ADR 0009：`{slot}` 既沒有起點也沒有間距，所以資料區不從第 1 列開始的
%           工作表定不到位。修法是 slot_origin + records（group_size / row_role）。
%        5. ADR 0010：兩個 row set 的等值合併。
%      這條線本身就是論文的論述：**表達力邊界是被資料集逼著往外推的，不是設計出來的**。
%
%      表達力邊界（寫進 §3.9，不要在這裡展開）：等值合併不做 outer join、不做範圍型
%      對應；聚合軸只有欄向；records 假設 group_size 固定；ADR 0008／0009／0010 的
%      驗證基礎都是 n=1。

\subsection{Failure Modes and What Each Mechanism Blocks}
\label{sec:design:loc:coverage}

% TODO 做成一張表，橫軸是第 2 章的失敗模式，縱軸是各機制。
%      這張表證明設計不是拼湊出來的。
% TODO 【現況要誠實寫】最近一次 derive 對映 31 條規則，21 條無法確認寫入位置，
%      成因是語意的（target_not_found / target_ambiguous）。這個數字要嘛先修掉，
%      要嘛作為 locator 機制現階段限制寫進 §3.9。不可以靜靜略過。

\section{Target Reference and Write-Back}
\label{sec:design:targets}

% TODO docx（bookmark / content control）、xlsx（具名範圍 / 座標）、dwg（handle）
%      三種目標參照的統一契約，以及各自的可寫入性判定。

\section{Deterministic Compiler and Zero-LLM Runtime}
\label{sec:design:runtime}

% TODO spec -> 可執行計畫；錯誤在編譯期而非執行期爆。受限 AST 求值。
% TODO value layer：實例輸入的驗證與型別。
%
% TODO ★ 論文級紅線，寫死在這一節 ★
%      執行期不得有 LLM fallback。related-work v1 審查判定 C4（zero-LLM execution）
%      是唯一乾淨的 claim，所以這不是實作偏好而是命題的必要條件。
%      具體推論：locator 定位失敗時必須失敗並回報，
%      不可以「叫 LLM 猜一個」來降低 §3.5.3 那個失敗率。

\section{Lifecycle Management}
\label{sec:design:lifecycle}

\subsection{Lineage and Versioning}
\label{sec:design:lineage}

% TODO spec 版本、產出版本、輸入版本三者的可追溯關係。

\subsection{Regression Suite}
\label{sec:design:regression}

% TODO 每次 spec 變更自動重跑全部案例。

\subsection{Drift Detection}
\label{sec:design:drift}

% TODO 模板改版時如何被發現、如何定位、如何最小化重新推導範圍。
%      fingerprint 對 docx dump 已驗證可用（文字改 -> 指紋不變、bookmark 改名 -> 指紋變）。
% TODO 【已知缺陷，修掉或寫出來】drift 目前只讀 target.handles 不讀通用的 targets，
%      一旦有 spec 使用通用 refs，漂移偵測會直接瞎掉。

\subsection{Auditability}
\label{sec:design:audit}

% TODO 為什麼這個架構能被工程流程接受，而 agent-per-instance 不能。
%      措辭注意：用 traceability / reproducibility / change control，
%      不要宣稱「法規要求 deterministic output」（找不到學術來源）。

\section{Expressiveness Boundary}
\label{sec:design:boundary}

% TODO 老實寫現階段表達不了的東西，這是第 6 章 future work 的來源：
%      - 寫入端不定長度明細（來源 7 筆就插 7 列）—— 插列會讓所有 positional ref 位移，
%        連動目標轉接器、比對、清空、fingerprint、drift。
%        替代形狀是「模板固定上限明細列位，未用列必須清空」，
%        而這正好把 negative obligation 變成一等公民
%      - 需要跨文件推理而非查表的欄位
%      - 【§3.5.3 的 locator 失敗率若沒修掉，也要寫進這裡】
 \newpage
% !TEX root = ../main.tex
\chapter{Methodology}
\label{ch:method}

% =====================================================================
% 骨架來源：docs/thesis/thesis-outline.md §5 —— 全篇最有研究味的一章
%
% 【本章與第 3 章的分界】
%   第 3 章 = 靜態存在的東西（規格長什麼樣、怎麼被執行）
%   第 4 章 = 它怎麼被產生出來的
% 對應樣板原本的 "System Model" -> "<你的演算法>" 結構，本章就是那個演算法。
% =====================================================================

\section{State Machine}
\label{sec:der:statemachine}

% TODO derive -> gate -> oracle -> repair -> converge / escalate to human

\section{Structural Gates}
\label{sec:der:gates}

% TODO schema 合法性、locator 可解析、覆蓋率門檻、內部一致性。
%      重點論述：這些在花錢跑實例之前就能擋掉大部分壞規格。
% TODO 【已量到的證據】一份壞規則不得摧毀產生它的 derive；
%      gate 擋下的東西要能被人審查而不是靜默丟棄。

\section{Two Complementary Oracles}
\label{sec:der:oracles}

% TODO ★ 本論文的核心實證洞見之一 ★
%      derive-time expected-text oracle：規格宣稱「這個欄位會產生什麼文字」，立即比對
%      case-level oracle：實際執行後與參考輸出比對
%      兩者抓到的錯誤類型不同，任何一方都無法取代另一方。
%      【要給證據與失敗案例分類，第 5 章有對應的消融實驗】
%
% TODO 為什麼不用 LLM-as-judge：在沉默失敗面前，judge 的錯誤與生成的錯誤相關。
%      這條論述同時支撐第 5 章「評測器不得詢問另一個 LLM」的決定。

\section{Case-Repair Loop}
\label{sec:der:repair}

% TODO 失敗案例如何回饋成修正指令；收斂性與停止條件；避免震盪（同一個 gate 反覆修）。
% TODO 反過擬合守則：修正僅限值層，不允許為了通過某個案例而改動結構層。

\section{Human in the Loop: Evidence Review}
\label{sec:der:evidence}

% TODO 人力從「每份文件都檢查」移轉為「一次性審查證據」。
%      要交代介面呈現什麼證據（欄位、locator 命中位置、預期值 vs 實際值）。

\section{Coverage Is the Real Bottleneck}
\label{sec:der:coverage}

% TODO ★ 核心實證洞見之二 ★
%      在達到高準確率之後，剩餘風險主要來自「規格沒涵蓋到的欄位」而不是「涵蓋了但算錯」。
%
% TODO 【現況與承諾的落差，必須誠實處理】
%      目前分母 100% 由 agent 自陳（rule_inventory），分子也是同一個 agent 的輸出。
%      後果實測得到：同一份實作跨兩輪 derive 報出 80% -> 83%「進步」，
%      原因只是第二輪把兩筆併成一筆。更根本的是：少枚舉一條規則分數就變高。
%
%      修法是 rule_segmenter（確定性切段，以內容雜湊為身分），讓規則文件成為權威分母。
%      實作未開始（TODO.md P0 FEATURE）。
%      若到寫作時仍未完成，本節必須把「規則文件是覆蓋率的權威來源」寫成規範性主張，
%      並在第 5 章 Threats to Validity 明列現行計分的偏誤方向。
%      【不可以照抄 deck 的說法】deck 承諾的是「可以對照解析後的規則文件實際計數」。
  \newpage
% !TEX root = ../main.tex
\chapter{Performance Evaluation}
\label{ch:evaluation}

% =====================================================================
% ★ 內容權威來源：docs/thesis/evaluation-protocol-v1.md ★
% 那份文件是完整的協定，本章是它的正文化。改動請先改那份，再同步過來。
% 舊的 thesis-outline.md §7.2–7.4 已作廢，不要回頭參照。
%
% 【2026-08-23 動筆註】§5.1／§5.2／§5.3／§5.5／§5.10 已寫成正文。體例沿用第 4 章：
% 每一節末尾一個 "Status:" 段落，主張要嘛 implementation-backed（指名檔案／實測數字），
% 要嘛 normative（明說尚未實證並揭露偏誤方向）。不允許沒有標記的中間態——第 3 章
% 骨架漂移的六處全部出在沒有標記的段落上。
% 尚未動筆的 §5.4／§5.6／§5.7／§5.8／§5.9 保留原本的 % TODO 寫作規格。
% =====================================================================

\section{Research Questions}
\label{sec:eval:rq}

The claim under test is not that a language model can fill in a form. It is that, for a
bounded class of recurring template-instantiation tasks, the model can be removed from the
recurring execution path altogether without losing correctness, traceability, or fail-loud
behaviour. Four research questions decompose that claim; a fifth asks which components of
the derivation procedure carry it.

\begin{description}
  \item[RQ1 --- Correctness.] Is a compiled specification's field-level semantic accuracy at
    least as good as a per-instance agent's on the same instances, measured by the same
    deterministic evaluator?
  \item[RQ2 --- Consistency.] How much does repeated execution of the same task vary? The
    question is asymmetric by construction, and the asymmetry is the point: for a
    per-instance agent the variance is present in every one of the $N$ executions, whereas
    for the compiled method execution is deterministic and \emph{all} of the variance sits in
    the single derivation. RQ2 therefore has two halves, and reporting only the first would
    be a misstatement: run-to-run variation of the baselines, and derivation-to-derivation
    variation of the compiled specifications.
  \item[RQ3 --- Cost structure.] What is the shape of cost against instance count? A one-off
    derivation cost plus a near-zero marginal execution cost crosses a per-instance cost at
    some $N$; the question is where that crossing lies, and how sensitive it is to derivation
    retries.
  \item[RQ4 --- Behaviour under template drift.] When the template changes underneath a
    frozen artifact, does the method fail loudly and locate the change, and do the baselines
    fail silently? This is the direct test of the design principle that governs the whole
    system.
  \item[RQ5 --- Component contribution.] Which of the structural gates, the two oracles, and
    the repair loop are load-bearing, and which are redundant with each other?
\end{description}

An earlier version of this plan carried a sixth question --- whether the method generalizes
beyond drawings to office documents. It has been removed rather than answered, because DOCX
and XLSX are now the primary evaluation carriers rather than a proof of concept, and a
question whose answer is a premise of the study design is not a research question.

\paragraph{A revision to the emphasis, and the measurement that forced it.}
The original framing put RQ1 first and treated RQ2 and RQ3 as secondary. A pilot measurement
on the DOCX dataset changed that ordering. Running a per-instance agent three times over one
calibration instance produced semantically correct output every time --- the cross-file
lookup, the rounding, the gram conversion, and the extent decision were all right --- and
three different files: three distinct output hashes in three runs, and cell-level agreement
of 119 of 172 cells, or 69.2\%. All 53 divergent cells were formatting decisions retaken from
scratch on each run, a thousands separator present in one run and absent in the next, and
none of them was a semantic error.

The consequence is stated here rather than in the results, because it is a claim about what
this thesis argues: the primary contribution is not that the compiled method is \emph{more
accurate} than a competent agent on tasks a competent agent can do. It is that accuracy is
comparable while consistency and cost structure differ structurally. Two tempting responses
to that pilot were considered and rejected. Raising task difficulty until the baseline fails
would tune the benchmark towards making the comparator fail; substituting a weaker model
would be choosing the opponent. Both are validity threats, and both are recorded as rejected
in Section~\ref{sec:eval:threats}.

\paragraph{Status: implementation-backed as a measurement, normative as a research plan.}
The pilot above was executed and its artifacts are retained. The five questions are a plan;
the formal experiment has not been run, and no figure in this chapter should be read as an
answer to any of them.

\section{Study Design}
\label{sec:eval:design}

\subsection{Why Not a Few-Shot Held-Out Split}
\label{sec:eval:whynotfewshot}

The conventional evaluation for a language-model system is a few-shot held-out split: the
model sees $K$ demonstrations and is scored on unseen tasks. That protocol borrows its
meaning from in-context learning, and this method is not that shape.

The template is an \emph{input}, not the object of generalization. The system is given the
template on purpose; so is every baseline, on every one of its executions. Seeing the
template is therefore not leakage but a premise, and it is symmetric across arms. This makes
the standard split either vacuous or unfair depending on how it is drawn. If the held-out
instances share a template with the demonstrations, nothing was held out that matters: the
compiled specification and the agent's demonstrations both describe that template. If the
held-out instances come from a different template, the compiled method must derive again ---
the demonstrations do it no good at all --- and the comparison measures the cost of a
derivation rather than generalization from examples.

What the split is genuinely needed for is narrower, and is retained: preventing the situation
in which one arm's prompts, gates, or intermediate representation were tuned on the very data
it is scored on while the baselines were not. The instrument for that is separation at the
\emph{template} level, and Section~\ref{sec:eval:split} reports how far that separation
actually holds in this study.

\paragraph{Status: normative, and deliberately so.} This subsection argues a design choice.
It is included because the objection it answers is the first one a reader of a
language-model-systems paper will raise, and because declining a standard protocol without
stating why reads as avoiding it.

\subsection{Two-Level Split}
\label{sec:eval:split}

The split has two levels, and only the first of them currently holds in full.

\paragraph{Instance level.}
Every evaluation template carries ten instances: $K=3$ calibration instances, shared by all
three arms and the only instances any arm may see during construction; two development
instances used for debugging; and $M=5$ sealed evaluation instances that no arm sees at any
stage. Two of the five evaluation instances in each dataset are cases whose \emph{correct}
behaviour is a refusal, so that an arm which always produces a file cannot score well merely
by producing files.

\paragraph{Template level.}
The evaluation templates are two, with a third planned:

\begin{itemize}
  \item \textbf{S4} (DOCX): a Commercial Invoice and a Packing List produced together from a
    product master file and a shipment detail file. Its principal difficulty is cross-file
    computation --- unit price from one workbook, quantity from another --- and it is the only
    DOCX evaluation dataset, so every DOCX claim in this thesis rests on it. Thirty-four
    applicable rules.
  \item \textbf{X1} (XLSX): the Directorate General of Highways monthly site-audit report
    chain, in which $N$ real 13020A audit sheets are aggregated into one 13020D monthly
    report. Its principal difficulties are a cross-file collection of source artifacts and a
    layout in which one logical record occupies two physical rows. Twenty-seven applicable
    rules.
  \item \textbf{S5} (XLSX, not yet built): a credit-programme completion audit, chosen for
    its density of negative obligations.
\end{itemize}

Two further items must be named, because their absence changes how the split reads.
\textbf{S3} is a synthetic pilot dataset that was built to protocol and then rejected on the
grounds that its layout and content were not those of a document that really circulates; it
is retained only as a pilot for the evaluator and the range locator, is not in the main
results table, and is not offered as evidence of external validity. \textbf{S1} and
\textbf{S2} appear in the protocol document as a proposed assignment of two development
templates and \emph{were never built}.

\paragraph{The gap this leaves, and the direction of its bias.}
Because S1 and S2 do not exist, development did not happen on templates disjoint from the
evaluation set. It happened on the DWG case study, on the rejected pilot, and --- decisively
--- on the calibration instances of S4 and X1 themselves. Three intermediate-representation
primitives were added specifically so that the compiled arm could produce X1 at all: artifact
collections, stacked records with a slot origin, and equality merge between resolved row sets.
X1 is an evaluation template. The template-level separation the protocol describes is
therefore, at present, a normative claim rather than a property of this study, and its bias
runs \emph{in favour of the proposed method}.

Following the third of the three admissible responses to such a gap, the claim is kept and
labelled rather than quietly weakened, and the bias is disclosed again in
Section~\ref{sec:eval:threats}. Two measurements bound its size without cancelling it. First,
each primitive was added only after it had been observed to fail on a real file, and its
collateral effect on the already-stored specifications was swept: the slot-alignment gate
introduced with the third primitive was replayed against all thirty-four stored
specifications and matched exactly the two it was written for, with no false positive on the
five specifications of the other dataset. Second, value-level verification of the resulting
specification was performed on development instances the derivation never saw: the fourth X1
specification reproduced 255 of 255 scored values across the three calibration and two
development instances using the scoring policy's own comparator, and a deliberately corrupted
variant of the same specification scored 240. Neither measurement makes the development
location correct; they bound how far it could have been exploited.

\paragraph{A second construct-validity item, recorded here and again in
Section~\ref{sec:eval:threats}.}
X1's rules require projects to be reported in ascending project-number order. The real 13020D
form prescribes no ordering; the ordering is a convention this dataset introduced so that the
extent has a checkable manifestation. The mechanism that implements it is built and tested,
but the question of \emph{who decided the order} belongs to dataset construct validity rather
than to the method.

\paragraph{Status: implementation-backed at the instance level, normative at the template
level, with the bias direction stated.}

\subsection{Dataset Construction}
\label{sec:eval:dataset}

% TODO 每個案例由人工撰寫的 generator 產生三樣：input.* / expected.* / ground\_truth.json
%      期望值必須從規則文件算出來，不得由受測系統產生——否則比對是循環論證。
% TODO 案例之間明細列數必須不同（3/7/15/40），且至少一個案例帶污染列。
%      否則外延規則根本沒被測到，一份寫死 D5:D19 的規格會全過。（ADR 0002 §4.1）
%
% TODO 標註規則已定案（2026-08-15），權威是 docs/thesis/annotation-manual-v1.md
%      ＋ 兩份 JSON Schema。這一節要寫的是它的兩層結構與為什麼是兩層：
%      每模板一份判準（manifest）、每實例一份值（ground\_truth）。
%      理由不是美觀而是一致性——24 個實例各自重複一次判準，任何一次抄漏都會讓
%      兩個實例用不同標準判分，而那種不一致在主結果表上看起來會像方法的差異。
% TODO 五種必要能力（協定原本列四種，標註定案時補上第五種）：
%      數值容差與格式、negative obligation、範圍外延（逐列列出而非只給起訖）、
%      semantic/presentation 雙層、以及 untouched（沒有任何欄位宣告的位置被動到）。
%      第五種的理由要寫：四種全滿足卻把模板其他地方清空的輸出仍可拿滿分。
% TODO 資料集約束再加一條：每個 aggregate 欄位必須配一個獨立的交叉驗證欄位，
%      理由見 §5.5——外延診斷只有 Arm C 宣告得出來，交叉驗證是唯一三個 arm 對等的偵測手段。

\section{Baselines and Fairness}
\label{sec:eval:baselines}

Three arms are compared. They differ only in what persists between instances.

\begin{description}
  \item[Arm A --- per-instance agent.] For each instance, an agent is given the template, the
    rule document, and the source data, and produces the output document. Nothing persists:
    the next instance starts from the same prompt with no artifact carried over. This is the
    status quo the thesis argues against, and it is deliberately a strong version of it ---
    a current frontier model with the same structured inspection tools the proposed method
    uses.
  \item[Arm B --- agent-built frozen skill.] The agent is first allowed to generalize from
    the calibration instances into a reusable artifact of its own design --- a script, a
    checklist, a set of instructions --- which is then frozen. Every evaluation instance is
    executed by an agent that may read that artifact and may not modify it.
  \item[Arm C --- the proposed method.] One derivation produces a specification; every
    instance is executed by the deterministic compiler with no model call at all.
\end{description}

\paragraph{Why three arms and not two.}
Arm B exists to answer the strongest available objection to the result. If only A and C were
compared, any advantage of C could be attributed to the mere existence of a persistent
artifact rather than to anything about how that artifact is produced and verified. Arm B
holds persistence fixed and varies only the production procedure: an agent generalizing
freely, versus a derivation constrained by an intermediate representation, structural gates,
and two oracles. The question Arm B answers is therefore the one that actually distinguishes
the contribution --- is the verified derivation better than an agent's own attempt at the
same idea?

\paragraph{Fairness has two layers, and the second is the one usually omitted.}
The first layer is \emph{tooling parity}: the same structured dump utilities, the same model,
the same retry ceiling, the same wall-clock ceiling. The second is \emph{material parity}:
Arms A and B receive an input list that is byte-identical to the one Arm C's derivation
receives, because an advantage produced by handing one arm a cleaner or richer set of inputs
is not an advantage of the method. The frozen input manifest is truth-free by construction and
its envelope is hashable, so what each arm was given is recorded rather than asserted.

Material parity has a specific consequence that must be stated because it works against the
proposed method. The rule material given to any arm may not contain target coordinates ---
no cell addresses, no bookmarks, no content-control tags, no drawing handles. Arm C's
derivation therefore has to locate its own write targets, and cannot be given them; the rule
table is not permitted to become an answer key for anyone.

\paragraph{Arm B's freeze must be enforced by the filesystem.}
An instruction not to modify the skill is not a control. The frozen artifact is hashed and
mounted read-only for the evaluation phase; an attempt to write to it is recorded as a
protocol violation and reported rather than silently tolerated. This matters because the
failure it prevents --- an agent quietly improving its own skill between instances --- would
convert Arm B into a per-instance arm while still being reported as a frozen one.

\paragraph{Status: contract implementation-backed, live execution not yet built.}
The three-arm input and evidence contract is frozen and tested: the truth-free input
manifest, the hashable envelope, the Arm B freeze, namespace isolation for Arm C, retention
of raw attempt evidence, and fake-capture negative controls are all implemented. The live
adapters that actually drive Arms A and B, and the formal wrapper around Arm C, are not.
The pilot reported in Section~\ref{sec:eval:rq} was run outside that harness and is labelled
as a pilot for exactly that reason; it must not be read as an Arm A result.

\section{Metrics}
\label{sec:eval:metrics}

% TODO 正確性：semantic 欄位級正確率（主指標）、案例級全對率（含 negative obligation）、
%              沉默失敗率、範圍外延正確率（診斷用）
%      presentation 分開報，不混算。不同載體不可以合併成單一正確率。
% TODO 可靠性：A/B 的 run-to-run 一致率與全 R 次皆正確的案例比例；
%              C 的執行期恆等（雜湊實證）+ 推導期變異度
% TODO 成本：derive 的 token/金額/wall-clock（一次性）、每實例成本、break-even N
% TODO 覆蓋率：以確定性切段後的規則文件為分母（現況見 §4.6）

\subsection{Asymmetric Rerun Budget}
\label{sec:eval:reruns}

% TODO ★ 要主動辯護這個設計 ★
%      A/B 各 R=5；C 主表 R=1 + 指定子集 R=5 以輸出雜湊實證恆等。
%      理由：C 的執行期是零 LLM 確定性的，跑五次必然逐位元相同，那不是實驗是複製檔案。
%
% TODO ★ 但 C 的變異度沒有消失，只是換了位置：它在推導期 ★
%      量法：同一模板獨立 derive D=3 次 -> 3 份規格 -> 每份跑全部評測實例，
%      報告三份規格之間的正確率／覆蓋率分佈與 spec diff 規模。
%      把 C 報成「變異度 0」而不交代這件事，會被審稿人一句話擊破。

\section{The Evaluator}
\label{sec:eval:evaluator}

All three arms are scored by one deterministic program. It parses the produced artifact,
extracts the values at the locations the dataset's annotation manifest names, compares them
against that instance's ground truth, and emits per-check results. It never asks a language
model whether an output looks right, for the reason given in Section~\ref{sec:der:oracles}:
a judge that can be wrong in the same direction as the system under test cannot bound the
system's error.

\subsection{Separate Channels, Never One Average}
\label{sec:eval:channels}

Scores are partitioned into channels, each with its own denominator: \emph{semantic} (the
primary metric), \emph{extent} (which source rows were included), \emph{refusal} (a case
whose correct behaviour is to refuse), \emph{structure} (the artifact opens and has the
expected shape), \emph{presentation} (formatting, reported separately from meaning),
\emph{untouched} (regions no field declares were not modified), and \emph{integrity} (the
package's critical parts survived). A check whose prefix is not one of the declared channels
fails closed rather than being scored as absent. The policy is versioned, and every result
record carries the evaluator version, the policy version, the source hash of the scoring
code, and the repository commit.

The partition is not presentational. The one measurement that settles it: while the ``this
slot must remain blank'' obligations were counted inside the semantic denominator, an arm
that submitted the blank template unchanged --- writing nothing whatsoever --- scored 66.7\%
semantic correctness on one instance. After the obligations were moved into their own
channel the same output scores 0\%. Every individual check had been correct; what was wrong
was which total they were added into.

\subsection{Two Self-Validations, Neither Sufficient Alone}
\label{sec:eval:selfvalidation}

\paragraph{Forward.} The evaluator must score the generator's own expected artifacts at
100\%. This catches an evaluator that reads the wrong location or compares with the wrong
tolerance.

\paragraph{Reverse.} Controlled corruptions are applied to the expected artifact and the
evaluator must catch every one. Without this, an evaluator that always answers ``pass'' would
give every arm a perfect score and the forward test would not notice. The current mutation
coverage is 6 of 6 on the pilot dataset, 12 of 12 on S4, and 13 of 13 on X1; X1 is the first
dataset in which \emph{all seven} channels have at least one mutation that exercises them,
which closed two channels that had previously been structurally vacuous.

\paragraph{A mutation test is not a validity proof, and the gap is specific.}
The expected artifact and the ground truth are produced by one generator. If the generator
computes a value wrongly, both are wrong together and the evaluator faithfully reports 100\%.
Mutation testing establishes that the evaluator detects corruption; it says nothing about
whether the generator computed correctly. The two do not overlap. The present detection
mechanism for that residual is a per-dataset spot audit --- at least three fields per dataset
recomputed by hand from the rule document, including one aggregate and one negative
obligation --- which is probabilistic rather than a guarantee, and is reported as such.

\subsection{A Negative Control for Every Zero}
\label{sec:eval:negativecontrols}

The evaluator and the probes around it are themselves software, and the failure mode observed
repeatedly during construction is that a broken probe reports a serene zero: zero mismatches,
zero problems, full marks. Measured instances during this work include a probe that used the
wrong key name, one that omitted a worksheet prefix from every reference so that nothing
matched anything, one whose trial compilation failed so that the comparison returned an empty
list, and one that compared string forms where the annotation declared exact decimal
comparison, so that twenty \emph{correct} cells were reported as wrong. Two of those report
``looks broken'' and the rest report ``looks fine''; the second kind is the dangerous one, and
it is indistinguishable from success by inspection.

The rule adopted is therefore mechanical rather than advisory. Any reported zero or perfect
score must be accompanied by a deliberate control that makes the figure move, and the control
must distinguish doing the work from not doing it --- a control whose expected value also
happens to be the result of skipping the mechanism measures nothing and never turns red. That
last clause is not hypothetical: an ordering test was found whose expected row order equalled
the arrival order, so disabling the sort entirely left it green.

\subsection{The Extent Diagnostic Is Not Symmetric Across Arms}
\label{sec:eval:extentasymmetry}

Only Arm C declares which source rows it included; Arms A and B produce a document and no
provenance. Scoring the extent channel from an arm's own declaration would repeat exactly the
mistake that the coverage denominator made in Section~\ref{sec:der:coverage} --- a figure
self-reported by the system being scored is not a cross-arm metric.

The cross-arm instrument is therefore a set of \emph{cross-checks}: identities between two
independently scored fields of the same instance, which any arm's output either satisfies or
does not. Alongside them, each dataset's manifest declares, for every wrong-extent hypothesis,
the values that hypothesis would produce, so a wrong extent can be \emph{named} from the
output alone rather than inferred. The effect is measurable: an output that occupies one extra
detail slot while leaving all six totals byte-identical passes all 112 semantic checks on one
instance and is still identified as a specific wrong-extent hypothesis. Arm C's own resolved
anchors are reported as supplementary evidence and are explicitly marked as available for one
arm only; they are never placed in a column alongside A and B.

\paragraph{Status: implementation-backed.}
The shared evaluator lints, self-tests, and scores all three datasets; the scoring policy is
frozen at version 1.1.0; the mutation suites and their per-channel negative controls are in
the repository and run in continuous integration. What is not established is the generator's
own arithmetic, as stated above.

\section{Template Drift}
\label{sec:eval:drift}

% TODO 漂移是獨立的一條軸，dev/eval 切分全是同模板，吃不掉它。
%      一個變體只動一個變因（挪動／改欄位名／換 block／增列擇一），否則失敗時分不出成因。
%      觀察：A/B 是否沉默地產生錯檔；C 是否在 gate 明確失敗、是否指出漂移位置、
%      重新推導需要多少成本。這是 fail-loud 設計原則的直接證據。

\section{Ablations}
\label{sec:eval:ablation}

% TODO 拿掉 structural gates／只留單一 oracle（證明 §4.3 的互補性）／拿掉 repair loop／
%      換模型（至少兩個，證明結論不綁單一供應商）。

\section{Results}
\label{sec:eval:results}

% TODO 實驗尚未執行。呈現形式先定：
%      Table 主結果（三 arm x 全部指標）
%      Table 消融
%      Fig. break-even 曲線（三條線）
%      Fig. 漂移情境對照
%      失敗案例定性分析——每一類失敗各給一個具體例子，這部分最能說服領域讀者
%
% TODO 主表的形狀已定（2026-08-15，評測協定 §5.3）：**逐模板一列，沒有單一總平均**。
%      跨載體合併本來就被禁止，所以主表是「每個評測模板一列 × 三個 arm」。
%      不要在正文寫出一個把 S3/S4/S5 平均起來的數字，那個數字沒有意義。
% TODO 外部來源模板（SpreadsheetBench 衍生，若採用）**獨立一小節，不進主表平均**。
%      理由要寫出來：S3/S4/S5 之間的差異是主困難不同（能力軸）；外部來源模板回答的是
%      「結論會不會只是因為 benchmark 是自己設計的」（效度探針）。把效度證據平均進
%      能力結論，兩邊都變模糊。它的 n 小是刻意的——要的是存在性陳述，不是檢定。
%      預期它的數字會比自建模板低（真實檔案較髒），而分開報才能把那句話寫成
%      「這一列是壓力測試」而不是被動解釋「為什麼被平均掉」。

\subsection{Where the Randomness Went}
\label{sec:eval:randomness}

% TODO 誠實的一小節，緊接在可靠性數字之後：
%      本方法沒有消除隨機性，只是把它從 O(N) 次移到 1 次，
%      並且在那 1 次上加了 gate 與 oracle。
%      這個交換為什麼值得——以及在什麼條件下不值得（N 太小、derive 成功率太低）。
%      後者的展開在 §6.2。

\section{Case Study: Engineering Drawings}
\label{sec:eval:dwg}

% TODO DWG 獨立報告，不併入主結果表。
%      價值在於證明方法在表徵最不友善、且交付物有法規與生產責任的載體上一樣成立。
%      現有規模：1 份基準模板、1 份規則文件、5 個回歸案例、4 個標準答案圖、0 個漂移變體。

\section{Threats to Validity}
\label{sec:eval:threats}

Each item below states what is threatened, what was measured, and which arm the bias favours.
Items whose direction favours the proposed method are listed first, because those are the ones
a reader has no independent way to discover.

\subsection{Threats Biased Towards the Proposed Method}
\label{sec:eval:threats:favourable}

\paragraph{T1 --- development happened on evaluation templates.}
The template-level separation described in Section~\ref{sec:eval:split} does not hold: the two
development templates the protocol proposes were never built, and the prompts, gates, and
intermediate representation were tuned against the calibration instances of the evaluation
templates themselves. Three representation primitives exist specifically because the compiled
arm could not otherwise produce X1. The bias favours Arm C, and it is not quantified --- the
only way to quantify it is to build the missing development templates. The two bounding
measurements are given in Section~\ref{sec:eval:split}.

\paragraph{T2 --- the cell-datatype exemption.}
The compiler stringifies every value before writing it, so a numeric field lands in a
spreadsheet as text. The annotation rules therefore declare \texttt{presentation.datatype} as
\texttt{any} for every field. Without that exemption Arm C would lose a point on every numeric
field, and the study would be measuring a write-side type policy rather than the method. With
it, a real defect --- numbers stored as text --- is not scored, and a baseline that writes
correct types earns nothing for doing so. This is the third admissible response to a gap: the
claim is kept, labelled normative, and the direction disclosed. It favours Arm C.

\paragraph{T3 --- the rule tables are ours.}
Even where the template is a genuine third-party artifact, the machine-readable rule
applicability manifest that supplies the coverage denominator was written for this study. It
weakens the objection that the templates are self-designed; it does not weaken the objection
that the rules were written by someone who knew what the method can express. Any use of an
externally sourced template must state this in the same breath, or it reads as padding.

\paragraph{T4 --- three primitives rest on a single example.}
Artifact collections, stacked records, and equality merge between row sets were each validated
on one dataset, and it is the same dataset in all three cases. The key--value extent added for
the DOCX dataset has the same property: a survey of six real government forms found no second
instance of the shape. The failure mode this predicts is not a red test but a different form
that silently does not fit, and no measurement in this thesis would detect it.

\subsection{Threats Biased Against the Proposed Method}
\label{sec:eval:threats:unfavourable}

\paragraph{T5 --- package preservation was, until recently, unmeasurable for Arm C.}
The integrity channel treats drawings and printer settings as critical parts of a spreadsheet
package, which is why the government form was chosen as a dataset in the first place. Until it
was corrected, the production write-back for XLSX serialized the workbook through a
spreadsheet model, and that operation drops exactly those two parts --- the same two parts the
mutation suite's integrity negative control drops, because the negative control \emph{is} a
round trip through the same library. Arm C therefore scored 0 of 5 on integrity by
construction, and the channel could not distinguish a correct run from a deliberately
corrupted one.

The response was the first admissible one: the write-back was replaced by a package editor
that edits the one worksheet part that must change and copies every other byte through. The
measurement after that change, on the same five instances, is 5 of 5 with no preservation
violation, while re-serializing the same outputs through the spreadsheet model still yields 0
of 5. Narrowing the channel to fit the writer was considered and rejected, because it would
have been a reduction of the thesis claim to match the implementation. The threat is recorded
rather than deleted, because every integrity figure produced before that correction is an
artifact of the writer and not a result.

\paragraph{T6 --- the DOCX write path has not had the same audit.}
The correction in T5 covers XLSX only. The layout-bearing parts of a Word package have not
been enumerated part by part, and S4 is the only DOCX evaluation dataset. Any DOCX integrity
figure in this thesis therefore carries an unmeasured residual.

\paragraph{T7 --- an unpublishable specification is not necessarily a failed one.}
The publish gate blocks on any unresolved item, and the unresolved set has repeatedly been
found to contain entries that are not defects. Two counters existed for the same quantity ---
the agent's own self-reported list and the gate's canonical set --- and they were never
compared: on the best X1 specification they read 13 and 47. Thirty-four of the 47 were detail
cells that a repeat block writes correctly, admitted to the gate's input list because the
author-time oracle reads expected text from that same list for a different purpose. Any count
of unresolved items must be decomposed before it is interpreted, and a headline
``unpublishable'' figure without that decomposition would understate the method.

\subsection{Threats of Undetermined or Mixed Direction}
\label{sec:eval:threats:mixed}

\paragraph{T8 --- dataset scale.}
Two evaluation templates exist and a third is planned, against an original intention of three
to five per carrier. Conclusions must be descriptive; no statistical significance may be
claimed, and no single average may be taken across carriers.

\paragraph{T9 --- the inclusion criterion is itself skewed.}
The datasets admit only tasks whose target row count is fixed, because the write side
deliberately does not insert rows, and they exclude formula results. Both exclusions coincide
with boundaries of the proposed method, so the excluded region may contain tasks the baselines
can do and the compiled method cannot. This is the second admissible response --- an explicit
expressiveness boundary --- and it is disclosed rather than argued away.

\paragraph{T10 --- what the untouched channel does not reach.}
Untouched-region comparison covers cells and paragraphs. It does not reach footnotes,
endnotes, comments, or other independent package parts, and styles and formula results are
outside the first version of the scoring policy by design, to avoid confounding correctness
with whether a spreadsheet application refreshed its cached values.

\paragraph{T11 --- model dependence.}
Every derivation reported here used one model family. The ablation in
Section~\ref{sec:eval:ablation} is the instrument for this and has not been run; until it is,
no claim in this thesis is established as vendor-independent.

\paragraph{T12 --- the correctness of the reference outputs.}
Discussed in Section~\ref{sec:eval:selfvalidation}. Mutation testing bounds the evaluator, not
the generator; the residual is covered only by a probabilistic spot audit.

\paragraph{T13 --- two responses to the pilot that were rejected, and why.}
After the pilot in Section~\ref{sec:eval:rq} showed the per-instance baseline to be
semantically accurate, two adjustments were available and both were declined. Increasing task
difficulty until the baseline begins to fail would tune the benchmark towards producing the
desired comparison, and substituting a weaker model would amount to selecting the opponent.
Recording them here is part of the claim: the emphasis of this thesis moved to consistency and
cost structure \emph{because} the baseline was accurate, not because the tasks were made
harder.

\paragraph{Status: implementation-backed for T1, T2, T5, T6, T7, T9, T10, T12; normative for
T3, T4, T8, T11, T13.}
The items marked implementation-backed each name a measurement in this repository. The others
are stated positions or unrun experiments, and are labelled so that no reader takes an
intention for a result.

% =====================================================================
% 【剩下的寫作規格】2026-08-23
% 已寫成正文：§5.1／§5.2（含三個小節）／§5.3／§5.5／§5.10。
% 仍是 % TODO 的：§5.4 Metrics（含 5.4.1 rerun budget）、§5.6 Drift、§5.7 Ablations、
%   §5.8 Results、§5.9 DWG case study。各節的註記就是它們的寫作規格，不要另開一份。
%
% TODO §5.3 的 Status 段承諾「A/B 的實作、預算、素材清單要完整交代」——
%      live adapter 建好之後要補一小節寫預算（重試上限、wall-clock 上限的實際數值），
%      現在寫會是憑空的數字。
% TODO §5.10 的 T5 與 §4.7 的 D4／D5 互相引用：三處講的是同一組修法的三個面向，
%      改任何一處都要同步。已於 2026-08-23 對齊過一次。
% TODO §5.8 Results 動筆時，§5.1 的 pilot 數字（69.2%）不可升格成 Arm A 的結果，
%      它不是走 harness 跑的。§5.3 的 Status 段已經寫了這句，兩處不可打架。
% =====================================================================
   \newpage
% !TEX root = ../main.tex
\chapter{Conclusions}
\label{ch:conclusion}

\section{Conclusion}
\label{sec:concl:summary}

% TODO 最後寫。回收 Introduction 的三個 claim，逐一說實驗結果支持到什麼程度。

\section{Limitations and When This Method Does Not Apply}
\label{sec:concl:limits}

% TODO 呼應 §3.1.5 的適用條件，但這裡談的是後果不是定義：
%      - 模板每次都不同（沒有 recurrence，過不了 break-even）
%      - 值域需要真實推理而非查表
%      - 一次性任務
%      - 規則無法以明確操作規則表達（例如需要專業判斷的裁量）
%      主動寫出來反而加分。與 §5.10 Threats to Validity 的分工：
%      那邊談「本實驗的結論可能怎麼不成立」，這裡談「本方法在什麼情境根本不該用」。

\section{Future Work}
\label{sec:concl:future}

% TODO 來源是 §3.9 的表達力邊界，不要另外發明：
%      - 寫入端不定長度明細（插列會讓所有 positional ref 位移，
%        連動目標轉接器、比對、清空、fingerprint、drift）
%      - 人工介入時間的量測（目前完全沒有量測工具）
%      - 更大規模的模板集與跨組織的規則文件形式差異
%      - 規則文件本身的品質問題（規則互相矛盾、規則缺漏）如何被偵測
   \newpage


% ====================================================
% 13. 參考文獻 Reference (檔案名稱 ref.bib)
%\ClearShipoutPicture % 把ref的浮水印關掉


\iftoggle{toc-use-cn} % 使用中文
{\addcontentsline{toc}{chapter}{參考文獻}
\renewcommand{\bibname}{參考文獻}} 
{\addcontentsline{toc}{chapter}{References}
\renewcommand{\bibname}{References}
} 


\bibliographystyle{IEEEtran}
\bibliography{ref}

}


% 14. 附錄 Appendix (用不到可以直接註解)
% 樣板原附的 appendix 是作者整理的 latex 排版提示, 不該當成本論文的附錄出版,
% 已移到 Sections/_samples/appendix.latex-tips.tex 保留參考。
% 之後若要放真的附錄 (例如完整的 spec IR schema、規則文件範例、失敗案例清單),
% 新建 Sections/appendix.tex 再把下面這行的註解拿掉。
%\input{Sections/appendix} \newpage


% =========================================
% [113年2月19日後上傳圖書館博碩士論文系統者適用] 
% 15. 學位論文發表形式確認書
% 16. 著作彙編之學位論文資訊及彙編學術著作之共同作者貢獻聲明書
% 這兩個不用附在論文PDF檔裡面

% 說明1: https://aa.nycu.edu.tw/userfiles/aach/files/20240125111443441.pdf
% 說明2: "國立陽明交通大學著作彙編之學位論文應行注意事項" https://bit.ly/457NTHk

% [心得感想]
% 畢業論文有使用到"就學期間發表"的論文, 就要填寫. 畢業後才刊登論文, 就屬於"不是".
% 對博士生來說, 在學期間可能發表很多篇論文, 可以放在畢業論文裡, 只是要記得填寫著作彙編的表格. 

% 著作彙編的說明
% https://ethics.moe.edu.tw/resource/epaper/html/27/


\balance
\end{document}